\section{Classi dello Schema dell'Ontologia}
\label{sec:tbox_classi}

\subsection{Gerarchia delle Classi}

L'ontologia definisce \textbf{12 classi top-level}, ciascuna corrispondente a
un'entit\`a fondamentale del dominio, pi\`u un ricco albero di sottoclassi che
raggiunge un totale di \textbf{78 classi}: 61 primitive, 15 definite tramite
\texttt{owl:equivalentClass} e 2 caratterizzate da un vincolo di cardinalit\`a massima
(condizione solo necessaria).

Le 12 classi top-level sono:
\begin{enumerate}[nosep]
  \item \texttt{vg:VideoGame} --- il videogioco, classe centrale dell'ontologia;
  \item \texttt{vg:Developer} --- lo studio di sviluppo;
  \item \texttt{vg:Publisher} --- l'editore;
  \item \texttt{vg:Genre} --- il genere videoludico;
  \item \texttt{vg:Platform} --- la piattaforma hardware/software;
  \item \texttt{vg:Character} --- il personaggio (giocante o non);
  \item \texttt{vg:Franchise} --- la serie/franchise;
  \item \texttt{vg:Award} --- il premio;
  \item \texttt{vg:GameEngine} --- il motore grafico;
  \item \texttt{vg:GameMode} --- la modalit\`a di gioco;
  \item \texttt{vg:Review} --- la recensione;
  \item \texttt{vg:GameEvent} --- l'evento del ciclo di vita.
\end{enumerate}

\subsection{Tabella Completa delle Classi Primitive}

La Tabella~\ref{tab:classi} elenca tutte le 61 classi primitive (escluse le 15
definite e le 2 con vincolo di cardinalit\`a, trattate nella Sezione~\ref{sec:classi_definite}), raggruppate per
classe top-level di appartenenza.

\begin{longtable}{llp{5cm}}
\caption{Classi primitive dell'ontologia}\label{tab:classi}\\
\toprule
\textbf{Classe} & \textbf{Parent} & \textbf{Descrizione} \\
\midrule
\endfirsthead
\multicolumn{3}{c}{\tablename\ \thetable{} -- continua dalla pagina precedente} \\
\toprule
\textbf{Classe} & \textbf{Parent} & \textbf{Descrizione} \\
\midrule
\endhead
\bottomrule
\multicolumn{3}{r}{\textit{Continua nella pagina successiva}} \\
\endfoot
\bottomrule
\endlastfoot

% ── VideoGame ──
\multicolumn{3}{c}{\textbf{VideoGame (classe top-level)}} \\
\midrule
VideoGame           & owl:Thing       & Prodotto videoludico \\
ReleasedGame        & VideoGame       & Gioco ufficialmente rilasciato \\
EarlyAccessGame     & VideoGame       & Gioco in accesso anticipato \\
CancelledGame       & VideoGame       & Gioco cancellato prima del rilascio \\
DelistedGame        & VideoGame       & Gioco rimosso dagli store digitali \\
IndieGame           & VideoGame       & Gioco a budget ridotto, sviluppatore indipendente \\
AAAGame             & VideoGame       & Gioco ad alto budget, grande publisher \\
GOTYWinner          & AwardWinningGame& Vincitore del premio Game of the Year \\
\midrule

% ── Developer ──
\multicolumn{3}{c}{\textbf{Developer (classe top-level)}} \\
\midrule
Developer           & owl:Thing       & Studio di sviluppo \\
FirstPartyDeveloper & Developer       & Studio proprietario di un platform holder \\
ThirdPartyDeveloper & Developer       & Studio indipendente \\
IndieDeveloper      & Developer       & Piccolo sviluppatore senza publisher \\
\midrule

% ── Publisher ──
\multicolumn{3}{c}{\textbf{Publisher (classe top-level)}} \\
\midrule
Publisher           & owl:Thing       & Casa editrice \\
FirstPartyPublisher & Publisher       & Divisione editoriale di un platform holder \\
ThirdPartyPublisher & Publisher       & Casa editrice indipendente \\
\midrule

% ── Genre ──
\multicolumn{3}{c}{\textbf{Genre (classe top-level)}} \\
\midrule
Genre               & owl:Thing       & Genere videoludico \\
ActionGenre         & Genre           & Azione, sfide fisiche \\
RPGGenre            & Genre           & Gioco di ruolo, progressione \\
StrategyGenre       & Genre           & Strategia, pensiero tattico \\
PuzzleGenre         & Genre           & Rompicapo, problem-solving \\
SimulationGenre     & Genre           & Simulazione di sistemi reali/fittizi \\
SportsGenre         & Genre           & Sport, competizioni atletiche \\
RacingGenre         & Genre           & Corse, competizione su veicoli \\
AdventureGenre      & Genre           & Esplorazione, narrativa \\
HorrorGenre         & Genre           & Horror, paura, suspense \\
PlatformerGenre     & Genre           & Piattaforme, ostacoli \\
\midrule

% ── Platform ──
\multicolumn{3}{c}{\textbf{Platform (classe top-level)}} \\
\midrule
Platform            & owl:Thing       & Piattaforma di esecuzione \\
ConsolePlatform     & Platform        & Console dedicata \\
PlayStationPlatform & ConsolePlatform & Famiglia Sony PlayStation \\
XboxPlatform        & ConsolePlatform & Famiglia Microsoft Xbox \\
NintendoPlatform    & ConsolePlatform & Famiglia Nintendo \\
PCPlatform          & Platform        & Personal computer \\
WindowsPlatform     & PCPlatform      & Microsoft Windows \\
MacOSPlatform       & PCPlatform      & Apple macOS \\
LinuxPlatform       & PCPlatform      & Linux \\
MobilePlatform      & Platform        & Dispositivo mobile \\
IOSPlatform         & MobilePlatform  & Apple iOS \\
AndroidPlatform     & MobilePlatform  & Android \\
\midrule

% ── Character ──
\multicolumn{3}{c}{\textbf{Character (classe top-level)}} \\
\midrule
Character           & owl:Thing       & Personaggio \\
PlayerCharacter     & Character       & Controllato dal giocatore \\
NonPlayerCharacter  & Character       & Controllato dall'IA \\
\midrule

% ── Franchise ──
\multicolumn{3}{c}{\textbf{Franchise (classe top-level)}} \\
\midrule
Franchise           & owl:Thing       & Serie/franchise \\
SubFranchise        & Franchise       & Spin-off o sotto-serie \\
\midrule

% ── Award ──
\multicolumn{3}{c}{\textbf{Award (classe top-level)}} \\
\midrule
Award               & owl:Thing       & Premio \\
IndustryAward       & Award           & Conferito da ente industriale \\
EditorialAward      & Award           & Conferito da testata editoriale \\
UserAward           & Award           & Determinato da voto popolare \\
\midrule

% ── GameEngine ──
\multicolumn{3}{c}{\textbf{GameEngine (classe top-level)}} \\
\midrule
GameEngine          & owl:Thing       & Motore grafico \\
ProprietaryEngine   & GameEngine      & Proprietario, non pubblico \\
OpenSourceEngine    & GameEngine      & Open-source \\
\midrule

% ── GameMode ──
\multicolumn{3}{c}{\textbf{GameMode (classe top-level)}} \\
\midrule
GameMode            & owl:Thing       & Modalit\`a di gioco \\
SinglePlayerMode    & GameMode        & Giocatore singolo \\
MultiplayerMode     & GameMode        & Multigiocatore \\
CoopMode            & GameMode        & Cooperativa \\
MMOMode             & GameMode        & Massively Multiplayer Online \\
\midrule

% ── Review ──
\multicolumn{3}{c}{\textbf{Review (classe top-level)}} \\
\midrule
Review              & owl:Thing       & Recensione (relazione N-aria gioco--revisore--punteggio) \\
\midrule

% ── GameEvent ──
\multicolumn{3}{c}{\textbf{GameEvent (classe top-level)}} \\
\midrule
GameEvent           & owl:Thing       & Evento del ciclo di vita \\
AnnouncementEvent   & GameEvent       & Annuncio pubblico \\
ReleaseEvent        & GameEvent       & Rilascio ufficiale \\
UpdateEvent         & GameEvent       & Aggiornamento/patch \\
DelistingEvent      & GameEvent       & Rimozione dagli store \\
\bottomrule
\end{longtable}

\subsection{Scelte di Modellazione}

\subsubsection{Partizione di Platform}

La classe \texttt{Platform} \`e stata partizionata in tre rami principali
(\texttt{ConsolePlatform}, \texttt{PCPlatform}, \texttt{MobilePlatform}),
ciascuno ulteriormente specializzato:

\begin{verbatim}
Platform
  +-- ConsolePlatform
  |     +-- PlayStationPlatform
  |     +-- XboxPlatform
  |     +-- NintendoPlatform
  +-- PCPlatform
  |     +-- WindowsPlatform
  |     +-- MacOSPlatform
  |     +-- LinuxPlatform
  +-- MobilePlatform
        +-- IOSPlatform
        +-- AndroidPlatform
\end{verbatim}

Questa struttura consente interrogazioni a diversi livelli di granularit\`a:
\textless{}giochi disponibili su console\textgreater{} (\texttt{ConsolePlatform}),
\textless{}giochi per PlayStation\textgreater{} (\texttt{PlayStationPlatform}),
oppure \textless{}giochi per PC Windows\textgreater{} (\texttt{WindowsPlatform}).
La partizione \`e resa semanticamente rigorosa da 4 blocchi \texttt{AllDisjointClasses}
(cfr.~Sezione~\ref{sec:disjointness}). Ogni piattaforma Wikidata viene classificata
al livello pi\`u specifico tramite le regole di reasoning (cfr.~Sezione~\ref{sec:reasoning}).

\subsubsection{Generi come Classi Aperte}

\texttt{Genre} ha 10 sottoclassi (\texttt{ActionGenre}, \texttt{RPGGenre}, \dots{})
ma \textbf{non} viene dichiarata come esaustiva (\texttt{owl:oneOf}) n\'e come
coperta da un'unione. Questa scelta \`e deliberata: Wikidata pu\`o contenere generi
non previsti dalla tassonomia attuale (es. \textless{}Battle Royale\textgreater{},
\textless{}Roguelike\textgreater{}). Un'ontologia chiusa impedirebbe l'aggiunta
di nuovi generi senza modificare lo schema. Con la modellazione aperta, eventuali
generi Wikidata non mappati restano istanze di \texttt{Genre} senza violare
la semantica.

\subsubsection{GameMode come Classe Top-Level Separata}

\texttt{GameMode} \`e modellata come classe top-level separata (anzich\'e come
attributo di \texttt{VideoGame}) per consentire il pattern \textbf{Value Partition}:
le sue 4 sottoclassi (\texttt{SinglePlayerMode}, \texttt{MultiplayerMode},
\texttt{CoopMode}, \texttt{MMOMode}) sono dichiarate mutuamente disgiunte,
garantendo che un gioco con pi\`u modalit\`a abbia asserzioni \texttt{hasGameMode}
distinte e semanticamente non ambigue. Questa scelta permette inoltre di definire
classi derivate (\texttt{SinglePlayerGame}, \texttt{MultiplayerGame}, \texttt{CoopGame})
tramite \texttt{owl:someValuesFrom} sulle modalit\`a.

\subsubsection{Review e GameEvent: Pattern N-ary Relation e Lifecycle Event}

\texttt{Review} e \texttt{GameEvent} sono state introdotte come nuove classi
top-level per modellare relazioni che richiedono reificazione:
\begin{itemize}[nosep]
  \item \textbf{Review} --- reifica la relazione tra un gioco, un revisore e un punteggio.
    Senza la classe \texttt{Review}, sarebbe impossibile associare un punteggio
    (datatype property) a una specifica combinazione gioco-recensore.
    Segue il pattern \textbf{N-ary Relation} (cfr.~Lez.~14);
  \item \textbf{GameEvent} --- reifica eventi temporali come annunci, rilasci e
    aggiornamenti. Consente di associare date e metadati a eventi specifici del
    ciclo di vita di un gioco, seguendo il pattern \textbf{Lifecycle Event}.
\end{itemize}
